\documentclass[10pt]{article}

\usepackage[utf8]{inputenc}

\usepackage[T1]{fontenc}

\usepackage{amsmath}

\usepackage{amsfonts}

\usepackage{amssymb}

\usepackage[version=4]{mhchem}

\usepackage{stmaryrd}

\usepackage{bbold}

\usepackage{graphicx}

\usepackage[export]{adjustbox}

\graphicspath{ {./images/} }



\begin{document}

намич. эквивалентности модельного и А. г. при произвольных температурах и для целого класса модельных систем (см. [3]). (О дальнейшем развитии и обобшении метода А. г. cM. [4]-[6].)



В качестве иллюстрации опишем нек-рые общие классы моделей, допускающих асимптотически точное решение (или существенное упрощение) методом А. Г., и приведем соответствующую каждому классу форму А. г. Везде будет подразумеваться, что гамильтониан $H$ (статистич. оператор $H=\Gamma-\mu N$, где $\Gamma$ - гамильтониан, $\mu-$ химич. потенциал, $N$ - оператор числа частиц) задан как самосопряженный оператор в подходящем гильбертовом пространстве и описывает систему частиц, заключенных в области $\Lambda$ конечного объема $|\Lambda|$ в $v$-мерном пространстве $\mathbb{R}^{v}$ или $\mathbb{Z}^{v}$. Введем обозначения: II. II - норма оператора, [·, .] - коммутатор, $A^{+}$- оператор, сопряженный $A, A^{\#}$ обозначает либо $A^{+}$, либо $\boldsymbol{A}$,



\[

f_{\Lambda}[H]=-(\beta / \Lambda)^{-1} \ln \operatorname{Tr} \exp (-\beta H)

\]



\begin{itemize}

  \item плотность термодинамич. потенциала и

\end{itemize}



\[

\langle A\rangle_{H}=\operatorname{Tr}\left(A e^{-\beta H}\right) / \operatorname{Tr} e^{-\beta H}

\]



\begin{itemize}

  \item среднее значение наблюдаемой $A$ в канонич. (большом канонич.) ансамбле Гиббса для системы с гамильтонианом (статистич. оператором) $H, \operatorname{Tr}(\cdot)-$ след оператора, $\beta=\left(K_{\mathrm{b}} T\right), K_{\mathrm{b}}-$ константа Больцмана, $T-$ абсолютная температура, $t$-lim - термодинамич. предел.

\end{itemize}



\begin{enumerate}

  \item Гамильтонианы с квадратичной по ограниченным интенсивным операторам формо й (см. [4], [5], [7]). Для простоты предполагается, что квадратичная форма дана в канонич. виде:

\end{enumerate}



\[

H=T-|\Lambda| \sum_{k=1}^{r} g_{k} A_{\Lambda, k} A_{\Lambda, k}^{+}+|\Lambda| \sum_{k=r+1}^{s} g_{k} A_{\Lambda, k^{A, k}} A_{\Lambda,}^{+},

\]



где $g_{k}>0, T=T^{+},\left|f_{\Lambda}[\mathrm{T}]\right| \leqslant M_{0}$ и для всех $k, k^{\prime}, k^{\prime \prime}=1, \ldots, s$ :



\[

\begin{gathered}

\left\|A_{\Lambda, k}^{\#}\right\| \leqslant M_{1}, \quad\left\|\left[A_{\Lambda, k}, A_{\Lambda, k^{\prime}}^{\#}\right]\right\| \leqslant M_{2}|\Lambda|^{-1}, \\

\left\|\left[A_{\Lambda, k}^{\#},\left[A_{\Lambda, k^{\prime}},\right]\right]\right\| \leqslant M_{2 T}|\Lambda|^{-1}, \\

\left\|\left[A_{\Lambda, k}^{\#},\left[A_{\Lambda, k^{\prime}}^{\#}, A_{\Lambda, k^{\prime \prime}}\right]\right]\right\| \leqslant M_{3}|\Lambda|^{-2},

\end{gathered}

\]



и, если $r<s$, то для всех $k=r+1, \ldots, s$ выполнено условие кластеризации



\begin{center}

\includegraphics[max width=\textwidth]{2023_07_08_7de28ee68f6d750e266dg-1}

\end{center}



где $\bar{b}_{\Lambda}(a)=\left\{\bar{b}_{\Lambda, k}(a)\right\}-$ решение системы уравнений $b_{k}=\left\langle A_{\Lambda, k}\right\rangle_{H_{0}(a, b)}$. В этом случае А. г. имеет вид



\[

\begin{gathered}

H_{0}(a, b)=T-|\Lambda| \sum_{k=1}^{r} g_{k}\left(a_{k} A_{\Lambda, k}^{+}+a_{k}^{*} A_{\Lambda, k}-a_{k} a_{k}^{*}\right)+ \\

+|\Lambda| \sum_{k=r+1}^{s} g_{k}\left(b_{k} A_{\Lambda, k}^{+}+b_{k}^{*} A_{\Lambda, k}-b_{k} b_{k}^{*}\right),

\end{gathered}

\]



причем (принцип минимакса):



\[

t-\lim \left|\inf _{a \in \mathbb{C}^{r}} \sup _{b \in \mathbb{C}^{(s-r)}} f_{\Lambda}\left[H_{0}(a, b)\right]-f_{\Lambda}[H]\right|=0 .

\]



Для существования термодинамич. предела для $f_{\Lambda}[H]$ необходимо потребовать существование $t-\lim f_{\Lambda}\left[H_{0}(a, b)\right]$ поточечно. Тогда



\[

\begin{gathered}

t-\lim f_{\Lambda}[H]=t-\lim \left\{\inf _{a \in \mathbb{C}^{r}} \sup _{b \in \mathbb{C}^{(s-r)}} f_{\Lambda}\left[H_{0}(a, b)\right]\right\}= \\

=\inf _{a \in \mathbb{C}^{r}} \sup _{b \in \mathbb{C}^{(s-r)}}\left\{t-\lim f_{\Lambda}\left[H_{0}(a, b)\right]\right\} .

\end{gathered}

\]



Примером системы с квадратичной формой общего вида является модель сверхпроводника с магнитным фазовым переходом, в к-рой учитывается как взаимодействие типа БКШ, так и редуцированное обменное взаимодействие коллективизированных электронов с локализованными магнитными моментами (см. [4]). Следует заметить, что в частном случае $r=s$, когда нет взаимодействия типа отталкивания, необходимость в условии кластеризации отпадает. К случаю $r=s=1$ относятся, напр., модели БКШ - Боголюбова и Кюри - Вейсса.



\begin{enumerate}

  \setcounter{enumi}{1}

  \item Гамильтонианы с функцией от ограни ченной, самосопряженной интенсивной на блюд а е мой (см. [4]): $H=T-|\Lambda| \varphi\left(A_{\Lambda}\right), \quad$ где $T=T^{+}$, $f_{\Lambda}[T]<\infty, \quad A_{\Lambda}=A_{\Lambda}^{+}, \quad\left\|A_{\Lambda}\right\| \leqslant M_{1}, \quad\left\|\left[T, A_{\Lambda}\right]\right\| \leqslant M_{2}$. Функция $\varphi(x)$ дважды дифференцируема на интервалах $x \in \mathbf{I} \equiv\left[-M_{1}\right.$, $\left.M_{1}\right]$, причем $\left|\varphi^{\prime \prime}(x)\right| \leqslant K, x \in \mathbb{I}$. Кроме того, для $a \in S_{\Lambda} \subset \mathbb{R}^{1}$ :

\end{enumerate}



\[

t-\lim \left\langle\left(A_{\Lambda}-a\right)^{2}\right\rangle_{H_{0}(a)}=0

\]



где А. г. имеет вид



\[

H_{0}(a)=T-|\Lambda| \varphi^{\prime}(a)\left(A_{\Lambda}-a\right)-|\Lambda| \varphi(a),

\]



a $S_{\Lambda}$ есть множество всех решений уравнения



\[

a=\left\langle A_{\Lambda}\right\rangle_{H_{0}(a)}

\]



Тогда



\[

t-\lim \left\{f_{\Lambda}[H]-\inf _{a \in S_{\Lambda}} f_{\Lambda}\left[H_{0}(a)\right]\right\}=0 .

\]



Если потребовать существование термодинамич. предела



для всех $x \in \mathbb{R}^{1}$, то существует



\[

F(x)=t-\lim f_{\Lambda}\left[T-x|\Lambda| A_{\Lambda}\right],

\]



\[

t-\lim f_{\Lambda}[H]=\inf _{a \in S}\left\{t-\lim f_{\Lambda}\left[H_{0}(a)\right]\right\}

\]



где $S \subset \mathbb{R}^{1}$ есть множество решений неравенств



\[

-F^{\prime}\left(\varphi^{\prime}(a)-0\right) \leqslant a \leqslant-F^{\prime}\left(\varphi^{\prime}(a)+0\right) .

\]



Эти неравенства впервые получены в [8] как обобщение уравнения самосогласования на случай, когда $F(x)$ не дифференцируема. Если $F(x)$ - дифференцируемая функция $x \in \mathbb{R}^{1}$, то



\[

S=\left\{a \in \mathbb{R}^{l}: a=t-\lim \left\langle A_{\Lambda}\right\rangle_{H_{0}(a)}\right\}

\]



Можно показать, что частный случай $\varphi^{\prime \prime}(a)>0, a \in \mathbb{I}$, описывает взаимодействие типа притяжения, а $\varphi^{\prime \prime}(a)<0, a \in \mathbb{I},-$ типа отталкивания. Нетривиальным примером системы этого класса служит модель, в к-рой $A_{\Lambda}$ есть гамильтониан плоской модели Изинга со взаимодействием ближайших соседей. Случай аналитич. функции многих переменных рассмотрен в [8].



\begin{enumerate}

  \setcounter{enumi}{2}

  \item Гамильтонианы вида (см. [4]):

\end{enumerate}



\[

H=H_{L}+\sum_{k=1}^{M} \omega_{k} b_{k}^{+} b_{k}+|\Lambda|^{-1 / 2} \sum_{k=1}^{M} \lambda_{k}\left(b_{k} L_{k}^{+}+b_{k}^{+} L_{k}\right),

\]



где $H_{L}=H_{L}^{+}, f_{\Lambda}\left[H_{L}\right]<\infty, b_{k}^{+}$и $b_{k}, k=1, \ldots, M,-$ бозе-операторы рождения и гибели, $\omega_{k} \geqslant \omega_{0}>0, \lambda_{k} \geqslant 0,\left\|L_{k}^{\#}\right\| \leqslant C|\Lambda|$. А. г. имеет вид



\[

\begin{gathered}

H_{0}(\eta)=H_{L}+|\Lambda| \sum_{k=1}^{M}\left\{\omega _ { k } ( | \Lambda | ^ { - 1 / 2 } b _ { k } ^ { + } + \lambda _ { k } \omega _ { k } ^ { - 1 } \eta _ { k } ^ { * } ) \left(|\Lambda|^{-1 / 2} b_{k}+\right.\right. \\

\left.\left.+\lambda_{k} \omega_{k}^{-1} \eta_{k}\right)-\lambda_{k}^{2}\left(\omega_{k}|\Lambda|\right)^{-1}\left(\eta_{k} L_{k}^{+}+\eta_{k}^{*} L_{k}\right)+\lambda_{k}^{2} \omega_{k}^{-1}\left|\eta_{k}\right|^{2}\right\},

\end{gathered}

\]



где $\eta=\left(\eta_{1}, \ldots, \eta_{M}\right) \in \mathbb{C}^{M}$. Тогда



\[

0 \leqslant \inf _{\eta \in \mathbb{C}^{M}} f_{\Lambda}\left[H_{0}(\eta)\right]-f_{\Lambda}[H] \leqslant O\left(|\Lambda|^{-1 / 3}\right),

\]



причем $O\left(|\Lambda|^{-1 / 3}\right)$ при $|\Lambda| \rightarrow \infty$ равномерно по температуре в любом ограниченном интервале.



Гамильтонианы этого типа используются для описания взаимодействия бозонного поля (фононы, фотоны) с веще-





\end{document}